\documentclass[11pt]{article}
\usepackage[margin=1in]{geometry}
\usepackage{amsmath,amssymb,amsthm,mathtools,longtable,array}
\usepackage[hidelinks]{hyperref}
\allowdisplaybreaks[2]
\setlength{\emergencystretch}{3em}
\newtheorem{theorem}{Theorem}[section]
\newtheorem{lemma}[theorem]{Lemma}
\newtheorem{proposition}[theorem]{Proposition}
\newtheorem{corollary}[theorem]{Corollary}
\theoremstyle{remark}\newtheorem{remark}[theorem]{Remark}
\newcommand{\E}{\mathbb E}
\newcommand{\dd}{\,d}
\newcommand{\Gp}{G_+}
\newcommand{\KK}{\mathcal K}
\newcommand{\BB}{\mathcal B}
\newcommand{\RR}{\mathcal R}
\newcommand{\eps}{\varepsilon}
\title{S2 closure package: gate budget, secondary energy,\\and retained weak-passage inputs}
\author{Analytical increment for review}
\date{Version 1; October 22, 2026}
\begin{document}\maketitle
\begin{abstract}
On the prescribed stopped interval, a temporary secondary-energy
guard pays the signed gate excess. Its coefficient enters only
vanishing errors. The resulting fixed global distance constant and
all-endpoint capped gain strictly improve that same guard through
the proved secondary lemma. TAILREP then pays the fixed strong
complement by mass. These are stopped-form original-flow results:
the fixed-chart c-form, core geometry and outside supremum remain
retained inputs. We compute the formerly unpaid secondary field
budget, and obtain a guard-independent strong learning-integral
budget from CFR's radial identity under the retained geometry.
The two Q3 strips and the complete c-form/passage assembly remain
open. No CLOCK passage theorem is applied.
\end{abstract}

\section{Setting, sources and retained guards}
All labels in this increment have prefix \texttt{an03s2cpv1:}.
Source abbreviations are those of roadmap v1.11: LCRC, GESG,
DECG, CFR, BULK, RCET, TAILREP, SEP, MOMC, PROFILE and CLOCK.
References to their named statements are external references;
none is reproved here. Exact paths and blob identifiers are in
the companion review note.

Use the initialized original flow, $P=(P_1+P_2)/2$,
$P_1=N(e_1,q^2I_2)$, $P_2=N(\rho e_2,q^2I_2)$, and set
\[
 \lambda=\rho^2\in[169/400,49/100],\quad Q=\lambda+q^2,
 \quad\ell=\log(1/q),\quad t_0=t_h.
\]
All cohorts are fixed label sets. The partition is
$\KK,\Gp,\Sigma,\BB$, with $\mathcal D=\varnothing$ and
$G=\Gp\cup\Sigma$. LCRC's balanced-coordinate convention gives
$E_S=H_S+A_S$, $\mu_S=E_S+J_S$. In particular
$H=H_{\mathcal C_q}\le H_+:=H_{\Gp}\le E_+:=E_{\Gp}$.
Write
\begin{equation}
 I(t)=Q\int_{t_0}^t(1-c),\qquad \beta=\lambda/Q,
 \quad r(t)=\sup_\theta|R_{11}(\theta,t)|,
 \quad b_0=Q(1-c_b).
 \label{an03s2cpv1:clock}
\end{equation}
The stopped guards on $[t_0,t_*]$ are
\begin{equation}
 \begin{gathered}
 0\le c\le c_b<1,\quad T_*:=t_*-t_0\le L_T\ell,\quad \int r\le B,\\
 J_G<q^2K_Y^2,\quad E_+<E_g,\quad
 \|e_2\|_{P_1}+\|e_1\|_{P_2}\le M_{\rm res}q.
 \end{gathered}
 \label{an03s2cpv1:guards}
\end{equation}
The diagonal loss bounds are the proved loss-monotonicity bounds.
Retain LCRC's instantaneous core chart $|\theta|,|\psi|\le Rq$,
mass bound $\mu_{\KK}\le M_K$, and the uniform outside budget
\begin{equation}
 b_B(q):=\sup_{t_0\le t\le t_*}\mu_{\BB}(t)/q=o(1).
 \label{an03s2cpv1:outsideguard}
\end{equation}
The constants $R,M_K$ are independent of the residual guards.
The proved S1 entry contracts are used with their actual reference:
MOMC for the right-half strong ancestry, Q2EC/Q2EP as incorporated
in SEP for $\Gp,\Sigma$, and the clock seed
$H(t_0)\asymp A_q^{-\lambda}q^k$, $k=10(1-\lambda)$.
Here $0<c_A\le A_q\le C_A$ are MOMC's fixed constants.
In particular $E_\Sigma(t_0)\le C_{\rm ent}q^5$.

\paragraph{The inherited local field guard.}
For the distance and gain conclusions, retain precisely the
DECG/GESG fixed-chart hypotheses: the same original reference is
in $[-R_c,R_c]^2$, $R_c=79/51$, with CLOCK offset $e<e_*$;
on $[-L_c,L_c]^2$, $L_c=R_c+1$, its receiver field has CLOCK's
c-form and its donor/sign hypotheses (in particular positive
donor measure and $M_s\le1$), with remainder
\begin{equation}
 \int_{t_0}^{t_*}\eps_F\le\overline F_F,
 \qquad \eps_F\text{ controls value and the induced }
 \ell^1\text{ receiver-derivative norm}.
 \label{an03s2cpv1:fieldguard}
\end{equation}
This is a retained original-flow field hypothesis, not a result
of this increment. The full budget $\int f_*$ from CLOCK is
\emph{not} assumed. The gate lemma below does not require
\eqref{an03s2cpv1:fieldguard}.

Fix $0<\nu\le1/100$, subsequently taking $q^2/\lambda\le\nu$.
Put
\begin{equation}
 \gamma=1-5\lambda+\frac{2\lambda}{1-\lambda},\qquad
 \Gamma_\nu=\gamma-\nu(k-2)
 \ge\frac{289169}{924000}>\frac3{10}.
 \label{an03s2cpv1:gamma}
\end{equation}
The inequality is BULK/LCRC's proved rational margin.

\section{The signed gate budget}
\begin{lemma}[A temporary coefficient occurs only in vanishing terms]
\label{an03s2cpv1:gate}
For any fixed $A<\infty$, supplement
\eqref{an03s2cpv1:guards}--\eqref{an03s2cpv1:outsideguard} by
\begin{equation}
 E_\Sigma(t)\le2A\{q^5+q^{\Gamma_\nu}H(t)\}.
 \label{an03s2cpv1:temporary}
\end{equation}
Choose $C_0<C_{\rm adm}(B,M_{\rm res})$ as in LCRC's localclock
lemma and then $q$ small. Uniformly over stopped endpoints,
\begin{align}
 \int_{t_0}^{t_*}(Qc+N_-)\dd t
 &\le \frac{10e}{9}\frac{E_g}{1-c_b}+o_A(1),
 \label{an03s2cpv1:gatebudget}\\
 \mathfrak X(t_0,t_*)&\le
 \frac{5e}{9}\frac{E_g}{1-c_b}+o_A(1),
 \label{an03s2cpv1:excess}
\end{align}
where $N_-=\E_{P_2}[(f_2X_2)_-]$ and
$\mathfrak X=\int\sup_\theta(|R_{22}|-I'/2)_+$.
The leading constants do not depend on $A$.
\end{lemma}
\begin{proof}
LCRC's localclock lemma uses the entry ratios and signs, response,
duration, diagonal loss bounds, cross residual and $\int r\le B$.
Its proof invokes neither $E_\Sigma$ nor any terminal amplitude
cap. It applies on this stopped interval and gives
\[
 z>0,\quad |a|\le z/3,\quad E_+\le(10/9)H_+,
 \quad H_+(s)\le eH_+(t)e^{-b_0(t-s)}.
\]
The primary defect is at most
$C(B,M_{\rm res})(C_0+q^2\ell)\le1$; the clock defect is
$D_{\rm clk}=O(q^{1/8}+q^2\ell)=o(1)$ separately. Hence
\begin{equation}
 \int E_+\le\frac{10eE_g}{9b_0},\qquad
 \int\sqrt{E_+}\le\frac{2\sqrt{10e/9}\sqrt{E_g}}{b_0}.
 \label{an03s2cpv1:primaryintegrals}
\end{equation}
No use of the secondary guard was made in these inequalities.

LCRC's exact response decomposition and responsecharge give
\[
 c=c_{\Gp}+c_\Sigma+c_{\KK}+c_{\BB},\qquad
 |c_S|\le E_S+\frac{2q}{\sqrt Q}\sqrt{E_SJ_S}.
\]
For the core and outside, pointcap's proof gives
$|c_{\KK}|\le Cq^2M_K$ and $|c_{\BB}|\le C\mu_{\BB}$.
Let $S_q=2A(q^5+q^{\Gamma_\nu}E_g)$. Then
$E_\Sigma\le S_q$, $J_\Sigma\le q^2K_Y^2$, and
\[
 \int Q|c_{\Gp}|\le\frac{10eE_g}{9(1-c_b)}
       +Cq^2K_Y\sqrt{E_g}/b_0,
 \qquad |c_\Sigma|\le S_q+Cq^2K_Y\sqrt{S_q}.
\]
GESG's proved equation \texttt{Ncohorts} is applicable because
the primary signs have just been obtained. It gives
\[
 N_-\le C\{q^2M_K+S_q+q^2K_Y\sqrt{S_q}
       +q b_B(q)+(E_g+q^2K_Y^2)e^{-k_1/q^2}\}.
\]
Thus an explicit upper bound on the vanishing error in
\eqref{an03s2cpv1:gatebudget}, after enlargement of a fixed $C$, is
\begin{align}
 \mathcal E_q(A)=C\bigg\{&\frac{q^2K_Y\sqrt{E_g}}{b_0}
 +L_T\ell\big[q^2M_K+S_q+q^2K_Y\sqrt{S_q}
 +q b_B(q)\nonumber\\[-2pt]
 &\hspace{42mm}+(E_g+q^2K_Y^2)e^{-k_1/q^2}\big]\bigg\}=o_A(1).
 \label{an03s2cpv1:error}
\end{align}
Finally GESG's signed lemma gives, without a pointwise cap on
$N_-$,
$\mathfrak X\le\tfrac12\int(Qc+N_-)
+\tfrac12M_{\rm res}q^2L_T\ell$.
This proves \eqref{an03s2cpv1:excess}.
\end{proof}

\begin{remark}[The factor $Q$ must be retained]
\label{an03s2cpv1:Qcorrection}
The direct consequence for the unweighted response is
\begin{equation}
 \int c\le \frac{10e}{9}\frac{E_g}{Q(1-c_b)}+o(1).
 \label{an03s2cpv1:cbudget}
\end{equation}
The same numerical expression without $1/Q$ does not follow
by dividing \eqref{an03s2cpv1:gatebudget}; $Q<1$ for small $q$.
This corrects the proposed side-benefit inference in the roadmap.
It leaves the needed $O(c_b)$ conclusion intact when $E_g=2c_b$,
since $Q\ge169/400$. No incorrect coefficient is used below.
\end{remark}

\section{The stopped S2 theorem and strict energy improvement}
Use continuous angle lifts and define the actual-reference distance
\[
 D_b(t)=q^{-1}(|\theta_b(t)-\theta_c(t)|+|\psi_b(t)-\psi_c(t)|),
 \qquad \xi_b=D_b(t_0).
\]
Thus $q\xi_b$ is the physical entry displacement in TAILREP's
gain-integration convention. For the strong tail use the fixed
right-half ancestry $\RR$.
Write $P=P_\rho=\Phi(\rho a_\rho)$,
$d=(1-\lambda P)/2$, $s=2d/(1-\lambda)$, $p=3/s$, and
$\alpha=\lambda P/2$. MOMC and PROFILE give, for
$\xi=D(t_0)$ and $\dd\eta_0=2m(t_0,b)\dd\lambda_0(b)$ on $\RR$,
\begin{equation}
 \eta_0\{\xi>u\}\le C\min\{1,(\epsilon_h/u)^p\},\qquad
 \epsilon_h=\sqrt2 A_q^d(1-M_0)^\alpha q^{10d-s},
 \quad\mu_{\RR}(t_0)=\tfrac12.
 \label{an03s2cpv1:strongentry}
\end{equation}
Here fixed multiplicative changes in $\epsilon_h$ are absorbed
in $C$, consistently with the normalized sandwich statements.
Take TAILREP's fixed-label core and complement
\begin{equation}
 \KK=\{b\in\RR:\xi_b\le d_c\},\quad
 d_c=c_*q^{\chi+kP/2},\quad \chi=1/20,\ c_*>0\text{ fixed},
 \quad \mathcal C=\RR\setminus\KK.
 \label{an03s2cpv1:cutoff}
\end{equation}
The future geometry of this core is still the retained state
input in Section 1; choosing this set does not establish it.

\begin{theorem}[S2 closure under the other stopped guards]
\label{an03s2cpv1:S2}
Under Section 1's retained guards, S1 entries and
\eqref{an03s2cpv1:cutoff}, there is a fixed $C_\Sigma$ such that
the temporary guard
\begin{equation}
 G_\Sigma:\quad E_\Sigma<2C_\Sigma(q^5+q^{\Gamma_\nu}H)
 \label{an03s2cpv1:jointguard}
\end{equation}
cannot be the first guard to fail. On the stopped interval the
following hold for all strong-complement and $\Sigma$ receivers.

Let $a_*=C_R(1+M_{\rm res})$ be GESG's cross-entry coefficient,
$a_0=\lambda(1-c_b)/2$, and $W_*=4R_c+a_*/a_0$. Set
\begin{align}
 p_*&=\Phi(-7L_c/10),\qquad
 C_T=2+p_*^{-1}+(7/10)\phi(0)/p_*,\nonumber\\
 C_{\rm loc}&=p_*^{-1}\exp(C_T\overline F_F),\nonumber\\
 C_D&=\sec(1/2)C_{\rm loc}
 \exp\left(B+\frac{5eE_g}{9(1-c_b)}+1\right)
 \left(1+\frac{W_*}{L_c-R_c}\right).
 \label{an03s2cpv1:CD}
\end{align}
For a fixed $0<\delta_0<1/4$, before the first physical exit
$qD=\delta_0$,
\begin{equation}
 D_b(t)\le C_D\xi_b e^{\beta I(t)/2}.
 \label{an03s2cpv1:distance}
\end{equation}
This includes initially exterior labels and later chart reentries.
At every endpoint, including after physical exit,
\begin{align}
 \left(\log\frac{m_b}{m_c}\right)'
 &\le I'\min\{1,A_0q^2\xi_b^2e^{\beta I}\}+b_q,
 &\int_{t_0}^t b_q&\le4B+o(1),\nonumber\\
 m_b(t)&\le C_m m_b(t_0)
       \max\{1,A_0q^2\xi_b^2e^{\beta I(t)}\}^{1+\nu},
 &A_0&=C_D^2\max(1,\delta_0^{-2}),\quad C_m=e^{6B+3}.
 \label{an03s2cpv1:gain}
\end{align}
Moreover
\begin{align}
 E_\Sigma(t)&\le C_\Sigma(q^5+q^{\Gamma_\nu}H(t)),
 \label{an03s2cpv1:energy}\\
 \sup_{[t_0,t_*]}\mu_{\mathcal C}/q&\le Cq^{\Delta_\nu},
 \qquad \int_{t_0}^{t_*}\mu_{\mathcal C}/q\le Cq^{\Delta_\nu}\ell=o(1),
 \label{an03s2cpv1:compmass}\\
 \Delta_\nu&=p(5(1-P)-s-\chi)-1
 -(1+\nu)[k(1-P)-2-2\chi]_+
 \ge\frac{17141311}{140024000}>\frac3{25}.
 \label{an03s2cpv1:Delta}
\end{align}
The theorem improves only $G_\Sigma$; it does not close the
field, core, reference, residual or full outside guards.
\end{theorem}
\begin{proof}
Apply Lemma \ref{an03s2cpv1:gate} with an arbitrary fixed $A$.
GESG's one-exit splice has coefficient
$\sec(1/2)C_{\rm loc}e^{B+\mathfrak X+\eps_q}
(1+W_*/(L_c-R_c))$, where $\eps_q=o(1)$ for fixed guards.
Choose $q$ last so that the combined vanishing errors are at
most one. This gives \eqref{an03s2cpv1:distance} with
\eqref{an03s2cpv1:CD}, independent of $A$.
Before physical exit the initially strong branch remains in
$|\theta|,|\psi|<1/2$ for small $q$, because the reference is
$O(q)$ and $\delta_0<1/4$. This checks GESG's sine-distance
comparability. The local c-form derivative is used only before
the first fixed-chart exit; the matrix-value estimate covers
all subsequent reentries. SEP's fixed entry cap need not fit
inside $L_c$.

DECG Corollary 3.3 now applies with its formerly missing distance
premise proved. It gives the first line of
\eqref{an03s2cpv1:gain} and the raw mass multiplier
$e^{6B+1/\beta+o(1)}\max(1,A_0q^2\xi^2e^{\beta I})^{1/\beta}$.
Since $1/\beta=1+q^2/\lambda\le1+\nu$, the displayed fixed
$C_m$ suffices for small $q$. This includes the radial reference
gain, so no additional radial mass factor is omitted.

The entry profile is SEP's proved raw radial submeasure, with
fixed cap $Z_\Sigma(K(C_0))$, total mass $O(q^3)$ and tail index
$(1-3\lambda/2)/d$. The reference, distance, gain, response,
duration and stopped clock hypotheses of LCRC's secondary lemma
are now all available. The lemma first gives
\[
 E_\Sigma\le C\{q^5+q^\gamma H+q^{\Gamma_\nu}H^{1+\nu}\}.
\]
Here $H\le E_+<E_g$ before any secondary improvement. Thus it
gives \eqref{an03s2cpv1:energy} with a constant $C_{\rm out}$
depending on $C_D,A_0,C_m,E_g$ and the fixed entry constants,
not on $A$. Fix
\begin{equation}
 C_\Sigma\ge\max\{1,C_{\rm out},C_{\rm ent}+1\},\qquad A=C_\Sigma,
 \label{an03s2cpv1:chooseSigma}
\end{equation}
and only now choose $q$ small in the gate lemma. At entry,
$E_\Sigma(t_0)<2C_\Sigma q^5$. For fixed $q$ the cohorts are
fixed, their energies are continuous, and
$E_\Sigma/(q^5+q^{\Gamma_\nu}H)$ has a strictly positive
denominator. At a proposed first equality $2C_\Sigma$, the
proved bound is at most $C_\Sigma$, a contradiction. This
is strict improvement without importing any other guard's
future conclusion.

For the strong complement, the stopped clock supplies
\[
 e^{I(t)}\le e^{D_{\rm clk}}H(t)/H(t_0)
       \le C E_g A_q^\lambda q^{-k}.
\]
Since $I\ge0$ and $\beta\le1$, $A_0q^2e^{\beta I}\le Cq^{2-k}$.
Apply TAILREP's tail-moment lemma to
\eqref{an03s2cpv1:strongentry} and the already proved
\eqref{an03s2cpv1:gain}; the raw entry mass is $\eta_0/2$.
It yields
\[
 \mu_{\mathcal C}\le
 C\epsilon_h^p d_c^{-p}
       \max\{1,Cq^{2-k}d_c^2\}^{1+\nu}.
\]
Writing $E=5(1-P)-s$ gives
$\epsilon_h/d_c\asymp q^{E-\chi}$ and
$q^{2-k}d_c^2\asymp q^{-[k(1-P)-2-2\chi]}$.
The two resulting powers are exactly
\eqref{an03s2cpv1:Delta}. TAILREP proves both its rational
lower bound and $p-2(1+\nu)\ge49/7550>0$; these are analytic
margins, not numerical premises. Integrating the uniform mass
bound over $T_*\le L_T\ell$ proves
\eqref{an03s2cpv1:compmass}. No second distance-growth factor
is charged after this mass calculation.
\end{proof}

\section{Constant order and the small energy guard}
\label{an03s2cpv1:order}
One valid acyclic order is
\begin{equation}
 \begin{gathered}
 (\text{model},R,M_K,B_s,c_b,\nu,\delta_0,c_*)
 \longrightarrow(B,E_g,K_Y,M_{\rm res})\\
 \longrightarrow(L_c,\overline F_F,C_D,A_0,C_m)
 \longrightarrow C_0
 \longrightarrow(K(C_0),Z_\Sigma,\text{entry constants})
 \longrightarrow C_\Sigma\longrightarrow q.
 \end{gathered}
 \label{an03s2cpv1:orderdisplay}
\end{equation}
Here $L_c=R_c+1$ may already be fixed at the first step.
The choice of $\overline F_F$ specifies a retained budget; it
does not prove that the original field meets it. $C_0$ satisfies
LCRC's original $C_0<C_{\rm adm}(B,M_{\rm res})$ and, when CFR's
primary derivative estimate is used, its receiver-boundary
margin for this $L_c$. Shrinking it is permitted.

The suggested order $C_\Sigma\to C_0$ cannot be read as
independence of $C_\Sigma$ from $C_0$: SEP's cap and tail constants
depend on $K(C_0)$. Moving $C_\Sigma$ after those constants fixes
this without a loop, because $C_D$ is independent of
$C_\Sigma$ and its appearance in \eqref{an03s2cpv1:error} is
vanishing. $M_{\rm res}$ dependence in $C_D$ is harmless in
the secondary and complementary mass consumers: it multiplies
$q^5,q^{\Gamma_\nu}H,q^{\Delta_\nu}$, not an order-one weak
energy. BULK's fixed threshold can be any
$0<\eta_*<(\delta_0/(2C_D))^2$.

The fixed factors $p_*^{-1}$ and $C_T$ enter through
$C_{\rm loc}=p_*^{-1}e^{C_T\overline F_F}$. S2 needs this to be
finite; it does not need $C_T\overline F_F=o(1)$. CLOCK has a
separate small forcing budget. Neither $R$ nor $M_K$, $B_s$ or
$c_b$ is enlarged in response to the secondary constant.

\paragraph{Using $E_g=2c_b$.}
LCRC's written proof uses the more conservative
$E_g^{\rm old}=2+2c_b$; it does not literally prescribe $2c_b$.
The latter choice has the required slack: after
\eqref{an03s2cpv1:energy}, its proved pointcap lemma gives
\begin{equation}
 H_+\le(9/8)c_b+o(1),\qquad E_+\le(5/4)c_b+o(1)<2c_b
 \label{an03s2cpv1:smallEg}
\end{equation}
for each fixed $c_b>0$ and then small $q$. All remaining charges
in that pointwise lemma vanish by the current guards and
\eqref{an03s2cpv1:energy}. One may retain the conservative
$B,K_Y,M_{\rm res}$ choices obtained with $E_g^{\rm old}$;
they are still upper bounds for the smaller energy guard.
The slack is $3c_b/4+o(1)$, not an arbitrary fixed $1/2$.
This verifies compatibility; it is not a new execution of
LCRC's joint residual bootstrap.

\section{Quick consequences for the retained-input audit}
\begin{corollary}[CFR's secondary obstruction is paid]
\label{an03s2cpv1:U}
On the interval of Theorem \ref{an03s2cpv1:S2},
\begin{align}
 \mathcal U_\Sigma
 :=\int(E_\Sigma+K_Y\sqrt{E_\Sigma})
 \le L_T\ell\big[&C_\Sigma(q^5+q^{\Gamma_\nu}E_g)\nonumber\\
 &+K_Y\sqrt{C_\Sigma}
       (q^{5/2}+q^{\Gamma_\nu/2}\sqrt{E_g})\big]=o(1).
 \label{an03s2cpv1:Ubound}
\end{align}
Thus CFR's stated secondary selected-moment and negative-part
charges, including its scaled receiver derivative, have
vanishing integrated cost. This conclusion does not assert
the complete field guard \eqref{an03s2cpv1:fieldguard}.
\end{corollary}
\begin{proof}
Use \eqref{an03s2cpv1:energy}, $H<E_g$ and
$\sqrt{x+y}\le\sqrt x+\sqrt y$. The positive margin in
\eqref{an03s2cpv1:gamma} pays the logarithm. CFR's displayed
secondary terms are bounded by this integrand, with additional
$q^2K_Y^2$ terms whose integrals also vanish.
\end{proof}

\begin{proposition}[The learning integral under retained core geometry]
\label{an03s2cpv1:Bs}
For the core \eqref{an03s2cpv1:cutoff}, under the same stopped
guards and retained core geometry,
\begin{equation}
 \int_{t_0}^{t_*}|1-\kappa_1|\dd t\le\log2+o(1).
 \label{an03s2cpv1:Bsbound}
\end{equation}
In particular a fixed $B_s=\log2+1$ is admissible, independently
of $(B,E_g,K_Y,M_{\rm res})$. This is not a proof of continued
core geometry.
\end{proposition}
\begin{proof}
CFR's proved radial state identity, for each core label, is
\[
 (\log m)'=1-\kappa_1+\delta_m,\qquad
 |\delta_m|\le C_{R,L_c}\{q^2(1+M_K+M_{\rm res})
                   +J_G+q\sqrt{E_GJ_G}+\mu_{\BB}\}.
\]
If needed take the fixed receiver radius in this identity to be
$\max(R,L_c)$; this does not enlarge the c-form guard's square.
Using $E_G\le E_g+o(1)$, its uniform integrated error is at most
\begin{equation}
 C\{q^2(1+M_K+M_{\rm res}+K_Y^2+K_Y\sqrt{E_g+1})L_T\ell
                +q b_B(q)L_T\ell\}=o(1).
 \label{an03s2cpv1:radialerror}
\end{equation}
Let $\mu=\mu_{\KK}$. Fixed chart geometry gives
$|\kappa_1-\mu|\le C_Rq^2\mu$, with vanishing integral.
Average the radial identity against the current core mass:
\[
 (\log\mu)'=1-\mu+\eta(t),\qquad
 d_q:=\int|\eta|=o(1).
\]
The entry tail \eqref{an03s2cpv1:strongentry} gives
$\mu(t_0)=1/2-O((\epsilon_h/d_c)^p)=1/2-o(1)$;
indeed $p(5(1-P)-s-\chi)>1$ by
\eqref{an03s2cpv1:Delta}. Put $w=\log\mu$, so $w(t_0)<0$.
The absolutely continuous positive-part identity yields
\[
 w_+(t)+\int_{t_0}^t(\mu-1)_+
       =\int_{t_0}^t\mathbf1_{\{\mu>1\}}\eta\le d_q.
\]
Therefore, using the integrated equation for $w$,
\[
 \int_{t_0}^t|1-\mu|
 =w(t)-w(t_0)-\int_{t_0}^t\eta
                   +2\int_{t_0}^t(\mu-1)_+
 \le-\log\mu(t_0)+4d_q.
\]
Adding $\int|\mu-\kappa_1|=o(1)$ proves
\eqref{an03s2cpv1:Bsbound}. All error coefficients may depend
on the fixed guards, but only multiply vanishing scales.
Thus $B_s$ may be fixed first, and $B=C_rB_s+2$ subsequently,
without circularity. If a previously chosen $B_s>\log2$ is
already in use, it need not be changed.
\end{proof}

\section{Retained-input table and exact remaining inequalities}
\label{an03s2cpv1:audit}
The following statuses refer to the stopped setting above,
not to an unconditional joint weak passage.

\begin{longtable}{>{\raggedright\arraybackslash}p{.18\textwidth}
 >{\raggedright\arraybackslash}p{.20\textwidth}
 >{\raggedright\arraybackslash}p{.54\textwidth}}
\textbf{Input}&\textbf{Status/source}&\textbf{Exact budget or remaining obligation}\\\hline
Secondary field charge&Paid here; CFR and Corollary \ref{an03s2cpv1:U}&
$\mathcal U_\Sigma=o(1)$, equation \eqref{an03s2cpv1:Ubound}.\\[4pt]
Fixed-chart c-form&Retained; CFR gives only pieces&
Full value and induced $\ell^1$ receiver-derivative remainder must satisfy
$\int\eps_F\le\overline F_F$. The exact norm is in
\eqref{an03s2cpv1:fullnorm} below.\\[4pt]
Strong-complement outside&Paid here; MOMC/PROFILE, TAILREP&
$\sup\mu_{\mathcal C}/q\le Cq^{\Delta_\nu}$ and
$\int\mu_{\mathcal C}/q\le Cq^{\Delta_\nu}\ell$.\\[4pt]
Original Q3 interior&Paid by GESG&
$\sup\mu_{\mathrm{int}}/q=O(q^9)$ and
$\int\mu_{\mathrm{int}}/q=O(q^9\ell)$.\\[4pt]
Two Q3 strips&Unpaid; GESG strip contracts&
Need both $\sup(\mu_{\mathcal S_\pi}+\mu_{\mathcal S_{3\pi/2}})/q=o(1)$
and their integrated counterpart, or a proved family reassignment
with the consumer contracts below.\\[4pt]
Core geometry/reference&Retained; CLOCK not applied&
Need $e<e_*$ and the core chart, with CLOCK's full forcing bound
$\int(c+|V|+K_w+T+\eps_F)\le\overline F_{\rm pass}
\le e_*/(4C_{\rm pass})$, as well as its core-entry conditions.\\[4pt]
Strong learning integral&Paid under retained geometry; CFR, Proposition \ref{an03s2cpv1:Bs}&
$\int|1-\kappa_1|\le\log2+o(1)$; a fixed $B_s>\log2$ works.\\
\end{longtable}

\subsection*{What has, and has not, been paid in the field}
CFR's primary selected term has integrated value/derivative cost
$C_{L_c}C_{\rm RCET}E_gC_0+o(1)$. Its secondary term now costs
$o(1)$ by \eqref{an03s2cpv1:Ubound}. The core second-output
charge is $O(q^2M_K\ell)$; the order-one core first output must
remain in the strong donor field. The complement and Q3 interior
outside charges are now integrable at PROFILE's $\mu/q$ cost.
The strip charges are not. In particular the remaining full
contract is genuinely
\begin{equation}
 \int_{t_0}^{t_*}\sup_{u\in[-L_c,L_c]^2}
 \left\{|F_{\rm orig}(u,t)-F_c(u,t)|_1
 +\|D_uF_{\rm orig}(u,t)-D_uF_c(u,t)\|_{1\to1}\right\}\dd t
 \le\overline F_F,
 \label{an03s2cpv1:fullnorm}
\end{equation}
with the specified donor representation and nonnegative $K_w$.
Naming non-core moments is not a proof of this entire estimate.
Any mass normalization, $M_s-1$ contribution and selected residual
left by that representation must be included. Proposition
\ref{an03s2cpv1:Bs} gives integrated core overshoot $o(1)$, not
the pointwise assertion $M_s\le1$ used by the local c-form lemma.
The latter is still part of the explicitly retained field contract.

\subsection*{Fixed ancestry decomposition of the outside population}
Up to null endpoints the initialized circle consists of the
right-half ancestry $\RR$, Q2 and Q3. Q2 is $\Gp\cup\Sigma$.
Consequently the present outside population is exactly
\[
 \BB=\mathcal C\ \dot\cup\ \mathcal Q_{3,\mathrm{int}}^0
          \ \dot\cup\ \mathcal S_\pi\ \dot\cup\ \mathcal S_{3\pi/2}.
\]
GESG takes $k_G=169/3200$,
$\Lambda_q=\sqrt{3\ell/k_G}$, $w_q=\arcsin(q\Lambda_q)$,
with initial interior $[\pi+2w_q,3\pi/2-2w_q]$ and the two
remaining strips of width $2w_q$. Its initialization-based
speed proof pays that fixed interior, not just an entry-time
moving mask. There is no other ancestry in this initialization.
The whole supremum \eqref{an03s2cpv1:outsideguard} remains a
guard until the strips are dealt with. Even that supremum
does not pay their time integral on an $O(\ell)$ interval.

For a $\pi$-strip extension, first classify its labels. A primary
part needs, at $t_0$, $z>0$, $|a|/z\le1/4$, $qj/z\le C_0$,
$J\le Cq^3$ and $E=o(1)$. A recruited part needs the raw measure
and actual-reference displacement to obey
\begin{equation}
 \begin{gathered}
 \xi\le Z_\pi,\quad \eta_\pi(\mathbb R_+)\le Cq^3,\quad
 \eta_\pi\{\xi>u\}\le Cq^3\min\{1,(h_A/u)^{p_2}\},\\
 h_A=A_q^dq^{10d-2s},\quad p_2=(1-3\lambda/2)/d.
 \end{gathered}
 \label{an03s2cpv1:pistrip}
\end{equation}
It is not legitimate to assign every label near $\pi$ to
the secondary family. A $3\pi/2$-strip strong extension needs
a bounded actual-reference entry chart and
\begin{equation}
 \eta_{\rm late}\{\xi>u\}\le C\min\{1,(\epsilon_h/u)^{3/s}\},
 \qquad \dd\eta_{\rm late}=m(t_0,b)\dd\lambda_0(b),
 \label{an03s2cpv1:latestrip}
\end{equation}
with the corresponding core cutoff or TAILREP tail moment.
Joining only after $t_0$ requires an arrival-matched version
with pre-arrival gain paid; a fictitious entry at $t_0$ is not
allowed. Any remainder still needs its own supremum and
integrated $\mu/q$ budget. Neither extension is proved here.

\subsection*{The remaining CLOCK forcing budget}
The $c$ term is now small: with $E_g=2c_b$,
\eqref{an03s2cpv1:cbudget} is at most
$20e c_b/[9(169/400)(1-c_b)]+o(1)$. This is a proved source
for that term; CLOCK's conclusion is not used as its premise.
For the candidate moments of CFR, with $O=\Gp\cup\Sigma\cup\BB$,
\[
 V_O=-\frac{\E_{P_2}[(f_O)_1X_2]}{qQ},\qquad
 K_O=\frac{\E_{P_1}[(f_O)_2X_1]}q,
\]
BULK's cancellation-free cross-output bound and
\eqref{an03s2cpv1:primaryintegrals} imply
\begin{equation}
 \int(|V_O|+(K_O)_+)
 \le \frac{C}{b_0}(E_g+K_Y\sqrt{E_g})
            +C\int\mu_{\BB}/q+o(1).
 \label{an03s2cpv1:VKbudget}
\end{equation}
Indeed each primary moment is bounded by
$C(E_++J_++\sqrt{E_+J_+}/q)$; integrate using
$J_+\le q^2K_Y^2$. The secondary part is bounded by
\eqref{an03s2cpv1:Ubound} and $q^2K_Y^2T_*$, and the outside
part costs $C\mu_{\BB}/q$. CFR also gives
\[
 \int(K_O)_-\le C\mathcal U_\Sigma
                   +C\int\mu_{\BB}/q+o(1).
\]
Thus choosing $K_w=(K_O)_+$ still requires paying the negative
part in the remainder. Equation \eqref{an03s2cpv1:VKbudget}
is a moment budget, not the missing full c-form identification.

For $T$, CLOCK/LOC consumes the first distance moment of the
strong complement in its strong donor measure. The new mass
bound pays the original angular field of that complement at
$\mu_{\mathcal C}/q$ cost. To use it as the literal $T$ term,
one must also identify that measure, its normalization and the
coordinates assigned to escaped labels. The exact consumer is
\begin{equation}
 \int_{t_0}^{t_*}T(t)\dd t\le F_T,
 \qquad T(t)=\int_{\KK^c}D\dd\eta_t
 \quad\hbox{in the chosen CLOCK donor representation}.
 \label{an03s2cpv1:Tremaining}
\end{equation}
Alternatively the complement can be charged to the field
remainder, but that requires explicitly constructing that
representation. We do not silently exchange these conventions
or evolve $T$.

The primary part of \eqref{an03s2cpv1:VKbudget} is
$O(c_b+K_Y\sqrt{c_b})$ with $E_g=2c_b$. Conservative LCRC
guard constants can be fixed uniformly for $c_b$ in a sufficiently
small fixed interval by using the larger old energy budget.
Thus $c$ and the primary candidate moments are compatible with
an arbitrarily small forcing allowance. But the exact outstanding
joint requirement is still
\begin{equation}
 \int(c+|V|+K_w+T+\eps_F)
 \le\overline F_{\rm pass}\le e_*/(4C_{\rm pass}),
 \label{an03s2cpv1:passremaining}
\end{equation}
including the strip budgets, \eqref{an03s2cpv1:fullnorm} and
the representation in \eqref{an03s2cpv1:Tremaining}.
This increment proves compatibility of the newly paid terms,
not \eqref{an03s2cpv1:passremaining}. It stops before applying
CLOCK, improving the core/reference guards, or running the
joint weak-passage assembly.
\end{document}
